% answers_short.tex —— 2011 年·四川卷·理科 · 简短答案速查
% 由 tools/make_short_answers.py 自 answers.tex 的 \gksolhead/\gkanswer 提取，
% 并与 2011/2011_数学真题_四川理_答案要点.md 逐题比对后生成，勿手改（改源 MD 后重跑）。

\begin{gknotice}
\textbf{使用说明}：本节只给最终答案便于快速核对，完整推导见解析版。答案为编写者独立推导，\textbf{非官方标准答案}。
\end{gknotice}

\gkpart{选择题}
\begin{center}\begin{tabular}{ccccccccccccc}
\toprule
题号 & 1 & 2 & 3 & 4 & 5 & 6 & 7 & 8 & 9 & 10 & 11 & 12 \\
\midrule
答案 & B & A & B & D & B & C & A & B & C & A & D & B \\
\bottomrule
\end{tabular}\end{center}

\gkpart{填空题}
\begin{itemize}[leftmargin=2.6em]
  \item[\textbf{13.}] $-20$。
  \item[\textbf{14.}] $16$。
  \item[\textbf{15.}] $2\pi R^2$。
  \item[\textbf{16.}] ②③。
\end{itemize}

\gkpart{解答题}
\begin{itemize}[leftmargin=2.6em]
  \item[\textbf{17.}] (1) 最小正周期 $2\pi$，最小值 $-2$；(2) 证明见解析
  \item[\textbf{18.}] (1) $\dfrac{5}{16}$；(2) 分布列见解析，$E\xi=\dfrac72$
  \item[\textbf{19.}] (1) 证明见解析（$CD=C_1D=\dfrac12$）；(2) $\dfrac23$；(3) $\dfrac13$
  \item[\textbf{20.}] (1) $a_1=d$，$a_2=d(1+d)$，$a_3=d(1+d)^2$；$d=-1$ 时不是等比数列，$d\neq-1$ 时为等比数列，公比 $1+d$；(2) $S_n=1+(nd-1)(1+d)^n$
  \item[\textbf{21.}] (1) $y=\sqrt2x+1$ 或 $y=-\sqrt2x+1$；(2) 定值为 $1$
  \item[\textbf{22.}] (1) 减区间 $\left(0,\dfrac{9}{16}\right)$，增区间 $\left(\dfrac{9}{16},+\infty\right)$，极小值 $\dfrac18$，无极大值；(2) 当 $a\leqslant1$ 时无解；当 $1<a\leqslant4$ 时 $x=3-\sqrt{5-a}$；当 $4<a<5$ 时 $x=3\pm\sqrt{5-a}$；当 $a=5$ 时 $x=3$；当 $a>5$ 时无解；(3) 大于 $\dfrac16$
\end{itemize}
